% LaTeX Template for AI Problem Analysis Output (v2.0)
% Strategic Analysis of AMC12A 2023 Problem 13
% IMPORTANT: Use LaTeX commands for all mathematical symbols, NOT unicode characters

\documentclass[12pt]{article}
\usepackage[utf8]{inputenc}
\usepackage{amsmath, amssymb, amsthm}
\usepackage{geometry}
\usepackage{xcolor}
\usepackage[most]{tcolorbox}
\usepackage{enumitem}
\usepackage{fancyhdr}

% Page setup
\geometry{margin=1in}
\setlength{\headheight}{14.5pt}
\pagestyle{fancy}
\fancyhf{}
\rhead{AMC12A 2023 Problem 13 Analysis}
\lhead{\today}

% Color definitions
\definecolor{problemconfig}{RGB}{230, 242, 255}    % Light blue
\definecolor{strategic}{RGB}{255, 230, 242}       % Light pink
\definecolor{tactical}{RGB}{242, 255, 230}        % Light green
\definecolor{techniques}{RGB}{255, 242, 230}      % Light yellow
\definecolor{verification}{RGB}{255, 230, 230}    % Light red
\definecolor{solution}{RGB}{230, 255, 255}        % Light cyan
\definecolor{competition}{RGB}{242, 242, 255}     % Light purple
\definecolor{proofwriting}{RGB}{240, 230, 255}    % Light lavender
\definecolor{gold}{RGB}{218, 165, 32}

% Box styles
\newtcolorbox{problemconfigbox}{
    colback=problemconfig,
    colframe=blue,
    title=Problem Configuration Analysis,
    fonttitle=\bfseries,
    arc=3pt,
    boxrule=1pt,
    breakable
}

\newtcolorbox{strategicbox}{
    colback=strategic,
    colframe=purple,
    title=Strategic Framework Application,
    fonttitle=\bfseries,
    arc=3pt,
    boxrule=1pt,
    breakable
}

\newtcolorbox{tacticalbox}{
    colback=tactical,
    colframe=green,
    title=Tactical Methodology Selection,
    fonttitle=\bfseries,
    arc=3pt,
    boxrule=1pt,
    breakable
}

\newtcolorbox{techniquesbox}{
    colback=techniques,
    colframe=orange,
    title=Conversion Techniques Application,
    fonttitle=\bfseries,
    arc=3pt,
    boxrule=1pt,
    breakable
}

\newtcolorbox{verificationbox}{
    colback=verification,
    colframe=red,
    title=Verification Strategy Design,
    fonttitle=\bfseries,
    arc=3pt,
    boxrule=1pt,
    breakable
}

\newtcolorbox{solutionbox}{
    colback=solution,
    colframe=cyan,
    title=Solution Roadmap,
    fonttitle=\bfseries,
    arc=3pt,
    boxrule=1pt,
    breakable
}

\newtcolorbox{competitionbox}{
    colback=competition,
    colframe=purple,
    title=Competition Strategy Integration,
    fonttitle=\bfseries,
    arc=3pt,
    boxrule=1pt,
    breakable
}

\newtcolorbox{proofwritingbox}{
    colback=proofwriting,
    colframe=purple!70,
    title=Proof-Writing Strategy,
    fonttitle=\bfseries,
    arc=3pt,
    boxrule=1pt,
    breakable
}

\newtcolorbox{keyinsightbox}{
    colback=yellow!20,
    colframe=gold,
    title=Key Insight,
    fonttitle=\bfseries,
    arc=3pt,
    boxrule=2pt,
    leftrule=5pt,
    breakable
}

\newtcolorbox{warningbox}{
    colback=red!10,
    colframe=red!50,
    title=Warning: Common Pitfall,
    fonttitle=\bfseries,
    arc=3pt,
    boxrule=2pt,
    leftrule=5pt,
    breakable
}

\newtcolorbox{tipbox}{
    colback=green!10,
    colframe=green!50,
    title=Pro Tip,
    fonttitle=\bfseries,
    arc=3pt,
    boxrule=2pt,
    leftrule=5pt,
    breakable
}

\begin{document}

\title{AMC12A 2023 Problem 13\\Strategic Analysis}
\author{Competition Math Framework v2.1}
\date{\today}
\maketitle

\section*{Problem Statement}

In a table tennis tournament, every participant played every other participant exactly once. Although there were twice as many right-handed players as left-handed players, the number of games won by left-handed players was 40\% more than the number of games won by right-handed players. (There were no ties and no ambidextrous players.) What is the total number of games played?

$$\textbf{(A) }15\qquad\textbf{(B) }36\qquad\textbf{(C) }45\qquad\textbf{(D) }48\qquad\textbf{(E) }66$$

\begin{problemconfigbox}
\subsection*{A. Problem Classification}
\begin{itemize}[leftmargin=*]
    \item \textbf{Problem Type}: To Find (number of games)
    \item \textbf{Competition Context}: AMC 12A 2023, Problem 13
    \item \textbf{Difficulty Band}: Mid-band (P13) - Moderate difficulty
    \item \textbf{Mathematical Domain}: Combinatorics/Logic (Tournament, Counting, Constraints)
    \item \textbf{Estimated Time}: 2-4 minutes
\end{itemize}

\subsection*{B. Problem Structure Identification}
\begin{itemize}[leftmargin=*]
    \item \textbf{Given Information}: 
    \begin{itemize}
        \item Round-robin tournament (everyone plays everyone once)
        \item Twice as many right-handed as left-handed players
        \item Left-handed won 40\% more games than right-handed
        \item No ties, no ambidextrous players
        \item Multiple choice with specific values
    \end{itemize}
    \item \textbf{Constraints}: 
    \begin{itemize}
        \item If $L$ = left-handed, $R$ = right-handed, then $R = 2L$
        \item Total players: $n = L + R = 3L$ (divisible by 3)
        \item Games won by left: $1.4 \times$ (games won by right)
        \item Total games = games won by left + games won by right
    \end{itemize}
    \item \textbf{Objective}: 
    \begin{itemize}
        \item Find total number of games played
        \item Must match one of the answer choices
    \end{itemize}
    \item \textbf{Key Observations}: 
    \begin{itemize}
        \item In round-robin, total games = $\binom{n}{2} = \frac{n(n-1)}{2}$
        \item This is a triangular number
        \item If right-handed won $r$ games, left-handed won $1.4r = \frac{7r}{5}$
        \item Total games = $r + \frac{7r}{5} = \frac{12r}{5}$
        \item Therefore total must be divisible by 12
        \item Must satisfy both: triangular number AND divisible by 12
    \end{itemize}
\end{itemize}

\subsection*{C. Strategic Orientation}
\begin{itemize}[leftmargin=*]
    \item \textbf{Problem Category}: Logic/Combinatorics with constraints
    \item \textbf{Solution Approach}: 
    \begin{itemize}
        \item Identify necessary conditions (triangular, divisible by 12)
        \item Eliminate answer choices
        \item Verify remaining choice(s)
    \end{itemize}
    \item \textbf{Cross-Domain Potential}: 
    \begin{itemize}
        \item Combinatorics: Counting, binomial coefficients
        \item Number Theory: Divisibility
        \item Logic: Constraint satisfaction
    \end{itemize}
\end{itemize}
\end{problemconfigbox}

\begin{strategicbox}
\subsection*{A. Psychological Strategy}
\begin{itemize}[leftmargin=*]
    \item \textbf{Confidence Calibration}: P13 - accessible with logical reasoning
    \item \textbf{First Impressions}: 
    \begin{itemize}
        \item Word problem requiring translation to mathematics
        \item Multiple constraints to track
        \item Answer choice elimination likely effective
    \end{itemize}
    \item \textbf{Emotional Management}: 
    \begin{itemize}
        \item Don't rush - carefully identify all constraints
        \item Write down key relationships
        \item Systematic elimination is powerful here
    \end{itemize}
\end{itemize}

\subsection*{B. Investigation Strategy}
\begin{itemize}[leftmargin=*]
    \item \textbf{Initial Exploration}: 
    \begin{itemize}
        \item Round-robin: $n$ players → $\frac{n(n-1)}{2}$ games
        \item This must be a triangular number
        \item Check which answers are triangular:
        \begin{itemize}
            \item (A) 15 = $\frac{5 \cdot 6}{2}$ $\checkmark$
            \item (B) 36 = $\frac{8 \cdot 9}{2}$ $\checkmark$
            \item (C) 45 = $\frac{9 \cdot 10}{2}$ $\checkmark$
            \item (D) 48: Would need $\frac{n(n-1)}{2} = 48 \Rightarrow n(n-1) = 96$, no integer solution $\times$
            \item (E) 66 = $\frac{11 \cdot 12}{2}$ $\checkmark$
        \end{itemize}
        \item Eliminate (D)
    \end{itemize}
    \item \textbf{Pattern Recognition}: 
    \begin{itemize}
        \item Games won: left = $1.4r$, right = $r$
        \item Total = $2.4r = \frac{12r}{5}$
        \item Must be integer, so total divisible by 12
        \item Check remaining choices:
        \begin{itemize}
            \item (A) 15: $15 = 12(1.25)$ $\times$
            \item (B) 36: $36 = 12(3)$ $\checkmark$
            \item (C) 45: $45 = 12(3.75)$ $\times$
            \item (E) 66: $66 = 12(5.5)$ $\times$
        \end{itemize}
        \item Only (B) 36 satisfies both conditions!
    \end{itemize}
\end{itemize}

\subsection*{C. Argument Construction Strategy}
\begin{itemize}[leftmargin=*]
    \item \textbf{Approach}: Constraint-based elimination
    \item \textbf{Key Steps}:
    \begin{enumerate}
        \item Establish that total games must be triangular
        \item Establish that total games must be divisible by 12
        \item Eliminate choices not satisfying these conditions
        \item Verify unique answer
    \end{enumerate}
    \item \textbf{Verification Points}: 
    \begin{itemize}
        \item Check triangular number condition
        \item Check divisibility by 12
        \item Verify consistency with player count divisible by 3
    \end{itemize}
\end{itemize}
\end{strategicbox}

\begin{tacticalbox}
\subsection*{A. Core Mathematical Tactics}
\begin{itemize}[leftmargin=*]
    \item \textbf{Combinatorial Counting}: 
    \begin{itemize}
        \item Round-robin tournament: $\binom{n}{2} = \frac{n(n-1)}{2}$ games
        \item Each pair plays exactly once
        \item Must be triangular number
    \end{itemize}
    \item \textbf{Ratio and Proportion}: 
    \begin{itemize}
        \item Left won 40\% more than right
        \item If right won $r$: left won $1.4r = \frac{7r}{5}$
        \item Ratio left:right = 7:5
        \item Total games = $\frac{12r}{5}$
    \end{itemize}
    \item \textbf{Divisibility Analysis}: 
    \begin{itemize}
        \item Total games must be divisible by 12
        \item Quick check on answer choices
        \item Powerful elimination tool
    \end{itemize}
\end{itemize}

\subsection*{B. Domain-Specific Tactics}
\begin{itemize}[leftmargin=*]
    \item \textbf{Combinatorics - Triangular Numbers}:
    \begin{itemize}
        \item Form: $T_n = \frac{n(n-1)}{2} = 1+2+\cdots+(n-1)$
        \item To check if $m$ is triangular: solve $\frac{n(n-1)}{2} = m$ for integer $n$
        \item Equivalently: $n^2 - n - 2m = 0$, check if discriminant $1+8m$ is perfect square
    \end{itemize}
    \item \textbf{Logic - Constraint Satisfaction}:
    \begin{itemize}
        \item Multiple necessary conditions
        \item Answer must satisfy ALL conditions
        \item Intersection of conditions often unique
    \end{itemize}
\end{itemize}

\subsection*{C. Conversion Technique Selection}
\begin{itemize}[leftmargin=*]
    \item \textbf{High-Probability Techniques}:
    \begin{itemize}
        \item Answer choice elimination (fastest)
        \item Constraint checking
    \end{itemize}
    \item \textbf{Medium-Probability Techniques}:
    \begin{itemize}
        \item Full algebraic setup
        \item Case analysis on number of players
    \end{itemize}
    \item \textbf{Verification}:
    \begin{itemize}
        \item Backwards check from answer
        \item Verify all conditions satisfied
    \end{itemize}
\end{itemize}
\end{tacticalbox}

\begin{techniquesbox}
\subsection*{A. Primary Technique: Constraint-Based Elimination}
\begin{itemize}[leftmargin=*]
    \item \textbf{Technique Name}: Multiple Necessary Conditions Elimination
    \item \textbf{When to Apply}: 
    \begin{itemize}
        \item When problem has multiple constraints
        \item When answer choices given
        \item When direct solution is complex
    \end{itemize}
    \item \textbf{Application - Complete Solution}: 
    
    \textbf{Step 1: Identify First Constraint - Triangular Number}
    
    In a round-robin tournament with $n$ players, total games = $\binom{n}{2} = \frac{n(n-1)}{2}$.
    
    This must equal one of the answer choices. Check which are triangular:
    \begin{align*}
    \text{(A) } 15 &= \frac{5 \cdot 6}{2} \quad \checkmark\\
    \text{(B) } 36 &= \frac{8 \cdot 9}{2} \quad \checkmark\\
    \text{(C) } 45 &= \frac{9 \cdot 10}{2} \quad \checkmark\\
    \text{(D) } 48 &: n(n-1) = 96 \text{ has no integer solution} \quad \times\\
    \text{(E) } 66 &= \frac{11 \cdot 12}{2} \quad \checkmark
    \end{align*}
    
    \textbf{Eliminate (D)}.
    
    \textbf{Step 2: Identify Second Constraint - Divisibility by 12}
    
    Let $r$ = games won by right-handed players.
    
    Then games won by left-handed = $1.4r = \frac{7r}{5}$.
    
    Total games = $r + \frac{7r}{5} = \frac{5r + 7r}{5} = \frac{12r}{5}$.
    
    For this to be an integer, the total number of games must be divisible by $\frac{12}{5}$, which means \textbf{divisible by 12} (since we need $\frac{12r}{5}$ to be an integer, $12r$ must be divisible by 5, so $r$ must be divisible by 5, making total = $\frac{12 \cdot 5k}{5} = 12k$).
    
    Check remaining choices:
    \begin{align*}
    \text{(A) } 15 &= 12 \cdot 1.25 \quad \times\\
    \text{(B) } 36 &= 12 \cdot 3 \quad \checkmark\\
    \text{(C) } 45 &= 12 \cdot 3.75 \quad \times\\
    \text{(E) } 66 &= 12 \cdot 5.5 \quad \times
    \end{align*}
    
    \textbf{Only (B) 36 satisfies both conditions!}
    
    \textbf{Step 3: Verify the Answer}
    
    With 36 games and $\frac{n(n-1)}{2} = 36$, we get $n(n-1) = 72 = 8 \cdot 9$, so $n = 9$ players.
    
    If 9 players total and twice as many right-handed as left-handed:
    \begin{itemize}
        \item Left-handed: 3 players
        \item Right-handed: 6 players
    \end{itemize}
    
    Check: $3 + 6 = 9$ $\checkmark$
    
    For games won: $\frac{12r}{5} = 36 \Rightarrow r = 15$ games won by right-handed.
    
    Games won by left-handed: $\frac{7 \cdot 15}{5} = 21$.
    
    Check: $21 = 1.4 \cdot 15$ $\checkmark$ and $21 + 15 = 36$ $\checkmark$
\end{itemize}

\subsection*{B. Alternative Techniques}
\begin{itemize}[leftmargin=*]
    \item \textbf{Method 2: Full Algebraic Setup}
    
    Let $L$ = number of left-handed players, $R$ = number of right-handed players.
    
    Given: $R = 2L$, so total players $n = 3L$.
    
    Total games: $\frac{n(n-1)}{2} = \frac{3L(3L-1)}{2}$.
    
    Let $g_L$ = games won by left-handed, $g_R$ = games won by right-handed.
    
    Given: $g_L = 1.4g_R$, and $g_L + g_R = \frac{3L(3L-1)}{2}$.
    
    Substituting:
    \begin{align*}
    2.4g_R &= \frac{3L(3L-1)}{2}\\
    g_R &= \frac{3L(3L-1)}{4.8} = \frac{5L(3L-1)}{8}
    \end{align*}
    
    For integer solution, numerator must be divisible by 8.
    
    Try small values of $L$:
    \begin{itemize}
        \item $L=1$: $n=3$, games = $\frac{3 \cdot 2}{2} = 3$ (not in choices)
        \item $L=2$: $n=6$, games = $\frac{6 \cdot 5}{2} = 15$ (choice A)
        \item $L=3$: $n=9$, games = $\frac{9 \cdot 8}{2} = 36$ (choice B)
        \item $L=4$: $n=12$, games = $\frac{12 \cdot 11}{2} = 66$ (choice E)
    \end{itemize}
    
    Check which gives integer $g_R$:
    \begin{itemize}
        \item $L=2$: $g_R = \frac{5 \cdot 2 \cdot 5}{8} = \frac{50}{8}$ not integer $\times$
        \item $L=3$: $g_R = \frac{5 \cdot 3 \cdot 8}{8} = 15$ $\checkmark$
        \item $L=4$: $g_R = \frac{5 \cdot 4 \cdot 11}{8} = \frac{220}{8}$ not integer $\times$
    \end{itemize}
    
    Answer: $L=3$, total games = 36.
    
    \item \textbf{Method 3: Additional Constraint - Players Divisible by 3}
    
    Since $n = 3L$ and we need triangular numbers:
    \begin{itemize}
        \item (A) 15: $n=5$ (not divisible by 3) $\times$
        \item (B) 36: $n=9$ (divisible by 3) $\checkmark$
        \item (C) 45: $n=10$ (not divisible by 3) $\times$
        \item (E) 66: $n=12$ (divisible by 3) $\checkmark$
    \end{itemize}
    
    Combined with divisibility by 12, only (B) works.
\end{itemize}

\subsection*{C. Technique Selection Justification}
\begin{itemize}[leftmargin=*]
    \item \textbf{Why Constraint Elimination is Best}: 
    \begin{itemize}
        \item Fastest method (under 2 minutes)
        \item Requires minimal calculation
        \item Multiple choice friendly
        \item Each constraint eliminates many choices
    \end{itemize}
    \item \textbf{Why Full Algebra is More Rigorous}:
    \begin{itemize}
        \item Proves uniqueness
        \item Finds all valid scenarios
        \item Good if answer choices not given
    \end{itemize}
    \item \textbf{Time Efficiency}: 2 minutes with elimination; 3-4 minutes with full algebra
\end{itemize}
\end{techniquesbox}

\begin{verificationbox}
\subsection*{A. Verification Level Selection}
\begin{itemize}[leftmargin=*]
    \item \textbf{Chosen Level}: Level 4 - Direct Substitution (30-40 seconds)
    \item \textbf{Justification}: 
    \begin{itemize}
        \item Mid-band problem with multiple constraints
        \item Worth verifying all conditions satisfied
        \item Quick to check backwards
    \end{itemize}
    \item \textbf{Time Budget}: 30 seconds for full verification
\end{itemize}

\subsection*{B. Verification Checklist}
\begin{itemize}[leftmargin=*]
    \item \textbf{Check triangular number}:
    \begin{itemize}
        \item $36 = \frac{9 \cdot 8}{2}$ $\checkmark$
        \item 9 players in tournament
    \end{itemize}
    \item \textbf{Check divisibility by 12}:
    \begin{itemize}
        \item $36 = 12 \cdot 3$ $\checkmark$
    \end{itemize}
    \item \textbf{Check player ratio}:
    \begin{itemize}
        \item 9 players = 3 left-handed + 6 right-handed
        \item Ratio: 6 = 2(3) $\checkmark$
    \end{itemize}
    \item \textbf{Check games won}:
    \begin{itemize}
        \item Right-handed won: $36/2.4 = 15$ games
        \item Left-handed won: $1.4 \cdot 15 = 21$ games
        \item Total: $15 + 21 = 36$ $\checkmark$
        \item Ratio: $21/15 = 1.4$ $\checkmark$
    \end{itemize}
    \item \textbf{Confirm answer choice}: Matches (B) 36 $\checkmark$
\end{itemize}

\subsection*{C. Alternative Verification Methods}
\begin{itemize}[leftmargin=*]
    \item \textbf{Method 1}: Check discriminant
    \begin{itemize}
        \item For triangular: $1 + 8(36) = 289 = 17^2$ $\checkmark$
        \item $n = \frac{1+17}{2} = 9$ $\checkmark$
    \end{itemize}
    \item \textbf{Method 2}: Verify no other choice works
    \begin{itemize}
        \item Checked: (D) not triangular, (A)(C)(E) not divisible by 12
        \item Only (B) satisfies all conditions
    \end{itemize}
    \item \textbf{Method 3}: Reasonableness
    \begin{itemize}
        \item 9 players is reasonable tournament size
        \item 3 left-handed vs 6 right-handed makes sense
        \item 21 vs 15 games won is plausible (left-handed win more despite fewer players)
    \end{itemize}
\end{itemize}
\end{verificationbox}

\begin{solutionbox}
\subsection*{A. Step-by-Step Solution}
\begin{enumerate}[leftmargin=*]
    \item \textbf{Identify Round-Robin Constraint}: 
    \begin{itemize}
        \item Total games = $\binom{n}{2} = \frac{n(n-1)}{2}$ (triangular number)
        \item Check answer choices: (D) 48 is not triangular $\times$
    \end{itemize}
    
    \item \textbf{Identify Divisibility Constraint}: 
    \begin{itemize}
        \item Games won: left = $1.4r$, right = $r$
        \item Total = $2.4r = \frac{12r}{5}$
        \item Must be divisible by 12
    \end{itemize}
    
    \item \textbf{Check Remaining Choices}: 
    \begin{align*}
    \text{(A) } 15 \div 12 &= 1.25 \quad \times\\
    \text{(B) } 36 \div 12 &= 3 \quad \checkmark\\
    \text{(C) } 45 \div 12 &= 3.75 \quad \times\\
    \text{(E) } 66 \div 12 &= 5.5 \quad \times
    \end{align*}
    
    \item \textbf{Unique Answer}: Only (B) 36 satisfies both conditions
    
    \item \textbf{Verify}: 
    \begin{itemize}
        \item 9 players: 3 left, 6 right (ratio 1:2 $\checkmark$)
        \item Games won: 21 left, 15 right (ratio 1.4:1 $\checkmark$)
        \item Total: 36 games $\checkmark$
    \end{itemize}
\end{enumerate}

\subsection*{B. Alternative Approaches Summary}
\begin{itemize}[leftmargin=*]
    \item \textbf{Constraint Elimination}: Fastest, uses triangular + divisibility
    \item \textbf{Full Algebra}: Set up equations for $L$, $R$, solve systematically
    \item \textbf{Case Analysis}: Try $L = 1, 2, 3, 4$ and check which gives valid answer
    \item \textbf{Additional Constraint}: Players must be divisible by 3
\end{itemize}

All methods yield (B) 36.
\end{solutionbox}

\begin{competitionbox}
\subsection*{A. Time Management}
\begin{itemize}[leftmargin=*]
    \item \textbf{Estimated Time}: 2-4 minutes total
    \begin{itemize}
        \item Problem reading and translation: 30 seconds
        \item Identify triangular constraint: 30 seconds
        \item Identify divisibility constraint: 45 seconds
        \item Check answer choices: 30 seconds
        \item Verification: 30 seconds
    \end{itemize}
    \item \textbf{Time Checkpoints}: 
    \begin{itemize}
        \item At 1 minute: Should have triangular constraint identified
        \item At 2 minutes: Should have divisibility by 12 established
        \item At 2.5 minutes: Should have answer (B) 36
    \end{itemize}
    \item \textbf{Priority Level}: High (P13 should be completed)
    \begin{itemize}
        \item Standard logic/combinatorics problem
        \item Answer choice elimination is powerful
        \item Good point value for time invested
    \end{itemize}
\end{itemize}

\subsection*{B. Risk Assessment}
\begin{itemize}[leftmargin=*]
    \item \textbf{Confidence Level}: 98\%+ after verification
    \begin{itemize}
        \item Two independent constraints
        \item Unique answer satisfying both
        \item Backwards verification confirms
        \item Clean logical reasoning
    \end{itemize}
    \item \textbf{Abandonment Criteria}: 
    \begin{itemize}
        \item Should not abandon - very accessible
        \item If stuck, use systematic answer checking
        \item Maximum 4 minutes before moving on
    \end{itemize}
    \item \textbf{Flag for Review}: 
    \begin{itemize}
        \item No, if both constraints clearly identified
        \item Yes, if any uncertainty about divisibility reasoning
    \end{itemize}
\end{itemize}

\subsection*{C. Answer Format}
\begin{itemize}[leftmargin=*]
    \item Multiple choice: Select (B) 36
    \item Double-check: $36 = \frac{9 \cdot 8}{2}$ $\checkmark$ and $36 = 12 \cdot 3$ $\checkmark$
    \item Bubble carefully on answer sheet
\end{itemize}
\end{competitionbox}

\begin{keyinsightbox}
\textbf{Key Insight}: This problem has two powerful necessary conditions that combine to uniquely determine the answer:

\textbf{Condition 1 - Triangular Number}:
$$\text{Total games} = \binom{n}{2} = \frac{n(n-1)}{2}$$

This eliminates (D) 48 immediately.

\textbf{Condition 2 - Divisibility by 12}:

If right-handed won $r$ games, left-handed won $1.4r = \frac{7r}{5}$ games.

Total games = $r + \frac{7r}{5} = \frac{12r}{5}$.

For this to be an integer, total games must be divisible by 12.

This eliminates (A), (C), and (E).

\textbf{Only (B) 36 satisfies both conditions!}

The beauty of this problem is that multiple choice makes it exceptionally fast—each constraint cuts the search space dramatically, and their intersection is unique. The rigorous algebraic solution confirms this is the only valid answer for tournament size up to reasonable bounds.

Verification: 9 players (3 left, 6 right), 36 games total, left won 21, right won 15. Perfect!
\end{keyinsightbox}

\begin{warningbox}
\textbf{Warning}: Common pitfalls:
\begin{enumerate}
    \item \textbf{Misunderstanding "40\% more"}: This means $1.4 \times$, not $0.4 \times$
    \item \textbf{Ratio confusion}: Left:Right games = 7:5, not 5:7
    \item \textbf{Triangular test}: Don't forget to check if $n(n-1) = 2 \times \text{answer}$ has integer solution
    \item \textbf{Divisibility error}: Total must be divisible by 12, not just by 5 or by 12/5
    \item \textbf{Not checking all constraints}: Need BOTH triangular AND divisible by 12
    \item \textbf{Player count}: Remember $n = 3L$ (total must be divisible by 3 for player ratio to work)
    \item \textbf{Game count vs player count}: Don't confuse the two
\end{enumerate}
\end{warningbox}

\begin{tipbox}
\textbf{Pro Tip}: Tournament problems and answer choice strategies

\textbf{Round-Robin Tournaments}:
\begin{itemize}
    \item Every pair plays once: $\binom{n}{2} = \frac{n(n-1)}{2}$ games
    \item Triangular numbers: $1, 3, 6, 10, 15, 21, 28, 36, 45, 55, 66, 78, 91, \ldots$
    \item Quick test: Is $1+8m$ a perfect square? If yes, $m$ is triangular
\end{itemize}

\textbf{Percent Increase Translation}:
\begin{itemize}
    \item "40\% more than" means multiply by 1.4 or $\frac{7}{5}$
    \item "50\% more than" means multiply by 1.5 or $\frac{3}{2}$
    \item "25\% more than" means multiply by 1.25 or $\frac{5}{4}$
\end{itemize}

\textbf{Constraint-Based Elimination Strategy}:
\begin{enumerate}
    \item Identify all necessary conditions from problem
    \item For each condition, eliminate answer choices
    \item Often only one choice survives
    \item Verify that choice satisfies all conditions
\end{enumerate}

\textbf{This Problem's Constraints}:
\begin{itemize}
    \item Must be triangular (from tournament structure)
    \item Must be divisible by 12 (from game ratio 7:5)
    \item Player count divisible by 3 (from player ratio 2:1)
\end{itemize}

\textbf{Competition Strategy}:
\begin{itemize}
    \item Word problems: translate carefully to mathematics
    \item Multiple choice: use elimination aggressively
    \item Two strong constraints often uniquely determine answer
    \item Verify answer by working backwards (takes 20-30 seconds)
    \item Budget 2-3 minutes for this type of problem
\end{itemize}

\textbf{Quick Triangular Check}:
For small numbers, memorize: 1, 3, 6, 10, 15, 21, 28, 36, 45, 55, 66...

Or use: $\Delta = 1 + 8m$ is perfect square $\Leftrightarrow$ $m$ is triangular.
\end{tipbox}

\section*{Final Answer}

Using constraint-based elimination:

\textbf{Constraint 1}: Round-robin tournament → triangular number → eliminates (D) 48

\textbf{Constraint 2}: Total = $\frac{12r}{5}$ where $r$ = games won by right-handed → divisible by 12 → eliminates (A), (C), (E)

Only (B) 36 satisfies both constraints.

\textbf{Verification}: 9 players (3 left, 6 right), 36 total games, left won 21, right won 15. All conditions satisfied.

$$\boxed{\textbf{(B) } 36}$$

\section*{Confidence Assessment}
\begin{itemize}[leftmargin=*]
    \item \textbf{Confidence Level}: 99\% 
    \begin{itemize}
        \item Two independent necessary conditions
        \item Only one answer choice satisfies both
        \item Triangular: $36 = \frac{9 \cdot 8}{2}$ $\checkmark$
        \item Divisible by 12: $36 = 12 \cdot 3$ $\checkmark$
        \item Backwards verification: 9 players (ratio 1:2), games won 21:15 (ratio 1.4:1) $\checkmark$
        \item Answer uniquely determined
    \end{itemize}
    \item \textbf{Verification Applied}: Level 4 - Direct Substitution
    \begin{itemize}
        \item Primary method: Constraint-based elimination
        \item Confirmed: Both triangular and divisibility constraints
        \item Backwards check: All conditions satisfied with $n=9$, $L=3$, $R=6$
        \item Consistency: $21 + 15 = 36$ and $21 = 1.4 \times 15$
    \end{itemize}
    \item \textbf{Flag for Review}: No
    \begin{itemize}
        \item Very high confidence
        \item Clean logical reasoning
        \item Unique answer from constraints
    \end{itemize}
    \item \textbf{Time Spent}: 
    \begin{itemize}
        \item Estimated: 2-3 minutes for complete solution
        \item Appropriate for P13
        \item Efficient use of answer choice elimination
    \end{itemize}
\end{itemize}

\end{document}

